\documentclass[conference,a4paper,10pt]{IEEEtran}

% Overleaf: Menu -> Compiler -> XeLaTeX
\usepackage[UTF8,fontset=fandol,scheme=plain]{ctex}
\usepackage{amsmath,amssymb}
\usepackage{booktabs}
\usepackage{graphicx}
\usepackage{enumitem}
\usepackage{xurl}
\usepackage{cite}
\usepackage{balance}
\usepackage{xcolor}
\usepackage{tikz}
\usetikzlibrary{arrows.meta,positioning,fit}
\usepackage[hidelinks,unicode]{hyperref}

\hypersetup{
  pdftitle={不完整证据下的可审计 APT 调查控制：信息边界、参数治理与证据获取},
  pdfauthor={匿名作者},
  pdfkeywords={APT 调查, 不完整证据, 主动取证, 信息边界, 参数治理}
}
\urlstyle{same}
\renewcommand{\abstractname}{摘要}
\renewcommand{\IEEEkeywordsname}{关键词}
\renewcommand{\figurename}{图}
\renewcommand{\tablename}{表}
\setlist{nosep}

\title{不完整证据下的可审计 APT 调查控制：\\
信息边界、参数治理与证据获取\\[-0.1em]
{\large Auditable APT Investigation Control under Incomplete Evidence:\\
Information Boundaries, Parameter Governance, and Evidence Acquisition}}

% 双盲投稿时保留匿名信息；终稿请替换为真实作者、单位和邮箱。
\author{\IEEEauthorblockN{匿名作者}
\IEEEauthorblockA{匿名单位\\anonymous@example.com}}

\begin{document}
\maketitle

\begin{abstract}
从审计日志构建、压缩和学习溯源图，以及把威胁情报对齐到本地事件，已经能够在既有证据表面上定位异常实体、恢复攻击路径并组织调查故事。然而，当日志只被部分保留、采集通道可能失败、不同来源的证据角色不一致且预算受限时，部分对齐结果仍不能回答下一步应获取什么证据、何时停止，以及当前证据最多允许支持何种结论。本文把这一困境形式化为信息边界约束的调查控制问题。系统以可回溯来源的 evidence claims 构造证据缺口状态；规划器只能读取预执行动作目标、版本化成本和公开状态，不能读取动作实际恢复集合、随机遮蔽、实现通道状态或结果标签；执行后再根据新增证据或零收益反馈更新状态，并以支持上限和 STOP 约束输出。

与仅固定一套人工参数报告结果不同，本文把成本、粒度阈值、证据 corroboration、收益—成本换算系数和通道先验纳入版本化治理。评估单位为 C07–C12 六个独立案例或攻击链；每个案例的 45 个 mask/intensity/seed 条件是案例内配对重复测量，合计 270 条条件，不能作为 270 次独立攻击。冻结 legacy 口径下，XGBoost、透明 M2 和 Logistic/M3b 分别在 270、270 和 267 条重复条件达到内部目标，平均成功成本为 3.8926、3.8704 和 4.0000。将全部动作统一为单位成本后，三者分别为 1.6852、1.7296 和 1.6630，且均为 270/270；AFA-VOI Myopic、AFA-VOI Rollout-H3 和 Depth-2 也保持全成功，但平均动作数分别为 1.7333、1.7778 和 1.8963。

在 40 个粒度阈值组合、6 个 k-of-n corroboration 配置和 7 个 M2 成本换算系数上，M2 的内部成功率保持 1.0 且不产生粒度越界，但最低动作序列一致率分别为 1.0000、0.8148 和 0.8407。先验敏感性进一步显示，将 planner prior 缩放为 0.75 时，AFA-VOI Myopic 出现 2 个成功损失，Rollout-H3 虽无成功损失但改变动作序列；Depth-2 出现 1 个损失与 1 个修复。1.25 倍缩放没有改变动作或结果。外部 actor/campaign 真值和分析师效用记录当前均为零，因此 actor accuracy 不可识别；Round 1 的 Claim 加权 $\kappa=-0.1455$，Intent Jaccard=0.3673、micro-F1=0.4878，旧粒度任务又因 A/B 源文件哈希相同而作废。本文因此不主张新的 actor classifier、跨组织泛化、SOTA 或任一规划器全局最优。其贡献是一个可执行、可审计且能暴露参数依赖性的调查控制接口，以及对其成立边界的可复核报告。

\end{abstract}

\begin{IEEEkeywords}
APT 溯源；不完整证据；主动取证；信息边界；参数治理；主动特征获取；可支撑结论
\end{IEEEkeywords}

\section{引言}
近年来，DEPCOMM、PROGRAPHER、KAIROS 和 MAGIC 等权威文献已经显著推进了审计图压缩、异常检测和攻击摘要恢复 \cite{ref1,ref2,ref3,ref4}的整体流程，而 POIROT、EXTRACTOR、CLIProv 和 APT-CGLP 等则聚焦于从查询图、不精确匹配、日志—情报对比学习和图—语言预训练缩小高层情报与低层事件之间的语义鸿沟 \cite{ref5,ref6,ref7,ref8}。在调研相关文献的过程中，真正未闭合的问题发生于“部分对齐之后、归因输出之前”。现场证据可能因保留窗口、采集代理、权限、网络遥测和数据模式差异而缺失；一个动作又可能因通道失效而没有任何产出。此时，没有匹配到某项行为不等于该行为不存在，自动相关图或 ATT\&CK 标签也不等于独立的 actor 或 campaign 真值。若动作的公开描述由其实际可恢复结果反向生成，离线规划器还会在执行前间接读到答案键。

针对以上问题，本文提出了一种 investigation control layer：输入是上游对齐结果、可回指来源的证据声明、候选采集动作、预算和结论上限；输出是下一条采集动作或 STOP。动作执行前，规划器只能看到版本化公开视图；执行后，系统才接收新增证据或零收益反馈并重新计算支持上限。语言模型可以作为离线语义编译器或解释器，但不是本文必要的在线决策构件。

这类系统的主要风险并不只来自模型误差。动作成本、粒度阈值、OR/AND 覆盖语义、收益权重和通道可靠性先验都可能成为“隐形结论开关”。因此，本稿把贡献组织为四点：

\begin{enumerate}[leftmargin=*,label=(\arabic*)]
  \item 定义部分对齐之后的证据获取控制任务，并把公开动作目标与隐藏实际恢复集合的分离落实为可测试的运行时接口。
  \item 给出可更新证据缺口状态、k-of-n corroboration、支持上限、通道反馈和显式 STOP 组成的闭环。
  \item 用版本化 profile、SHA-256、正式运行 Gate 和配对敏感性分析治理成本、阈值、奖励和先验，而不是用一次重新调参替代依据。
  \item 在六个独立案例上同时报告正结果、排序反转、先验敏感性、人工构念失败和外部真值缺口，从而限定可声称与不可声称的结论。
\end{enumerate}

\section{相关工作与任务边界}
\subsection{溯源图、CTI 对齐与归因推理}
现有上游工作主要回答怎样从日志恢复攻击相关结构、怎样把报告行为映射到本地证据，以及怎样基于已有证据组织归因推理。它们的输出可以成为本文的 evidence claims 或候选行为节点，但不自动给出下一项取证控制指令。近年的真实长度 CTI 评估还显示，生成式模型在准确性、一致性和校准方面仍有明显风险 \cite{ref9}；高层 TTP 对高风险行为体也未必具有足够区分力 \cite{ref10}。因此，本文不把语言完整性、相似度或内部 coverage 直接解释为归因正确性。

\subsection{主动特征获取与数字取证规划}
主动特征获取（AFA）把未观测特征视为有成本动作，在预测收益与采集代价之间进行序贯权衡 \cite{ref11}。AACO、NOCTA 和 AFABench 分别提供非贪心采集、计划级 objective–cost 权衡和统一基准接口；WinRegRL 则在专家指定的 Windows 取证 MDP 中使用动态规划和强化学习 \cite{ref12,ref13,ref14,ref15}。这些工作已经覆盖“一般意义上有成本地选择下一观测”，故 AFA、MDP 或非短视搜索本身不是本文的新颖点。

本文的差异在部署信息结构：安全动作可能失败，动作可能返回一组证据而不是揭示一个静态特征，预执行目标不能等同于实际恢复集合，来源投影与原始事件需要区分，最终输出还受证据支持上限约束。本文的 AFA-VOI、Rollout-H3 和 Depth-2 都是同一 Project05 端点下的领域适配；AFABench、AACO、NOCTA 和 WinRegRL 的官方任务与本端点不等价，因而不称为官方同任务复现。

\subsection{成本与证据治理}
采集成本在成本敏感学习中是外生但必须显式的决策量 \cite{ref16,ref17}。信息价值理论要求把信息收益和代价映射到同一效用轴 \cite{ref18}。数字取证实践则以易失性次序和采集努力安排证据获取 \cite{ref19,ref20}。这些依据能够支持“为什么要分解和测量成本”，却不能给出本项目唯一正确的 72 维动作成本向量。因此本文不宣称找到真实成本，而是把 legacy、uniform、rubric 和 measured 作为不同状态的 profile：前两者已有正式运行，后两者必须在评分或测量完成并冻结后才能运行。

\section{问题定义}
\subsection{状态、动作和统计单位}
设攻击行为图为
\[
G=(V,E),
\]
时刻 \(t\) 可见且可回指来源的证据声明为 \(C_t\)。调查状态定义为

\[
s_t=(V_t^{\mathrm{cov}},E_t^{\mathrm{cov}},U_t,K_t,D_t,g_t,B_t,H_t),
\]

其中 \(V_t^{\mathrm{cov}}\) 和 \(E_t^{\mathrm{cov}}\) 分别表示已覆盖节点和边，\(U_t\) 表示未覆盖节点，\(K_t\) 表示关键缺口，\(D_t\) 表示阶段与证据类型摘要，\(g_t\) 表示当前可支撑结论粒度，\(B_t\) 表示剩余预算，\(H_t\) 表示动作和反馈历史。

候选动作 \(a\) 具有公开目标 \(I(a)\)、采集通道 \(q(a)\)、版本化成本 \(c_p(a)\) 和可选的公开先验；实际可恢复声明集合 \(R(a)\) 对非 Oracle 规划器隐藏。STOP 定义为零成本公共动作：

\[
c_p(a_{\mathrm{STOP}})=0.
\]

实验中的独立分析单位为 case 或 attack chain；mask strategy、mask intensity 与 seed 仅作为同一案例内的重复测量变量。本文主评估包含 \(6\) 个独立案例和 \(270\) 个配对重复条件，因此：

\[
N_{\mathrm{case}}=6,\qquad N_{\mathrm{repeat}}=270 .
\]

其中 \(270\) 条条件不作为独立样本，不进行将重复条件视为独立观测的显著性检验，即：

\[
n_{\mathrm{ind}}=6\neq270 .
\]

\subsection{k-of-n corroboration 与支持上限}

对节点 \(v\) 的 required claims 集合记为 \(Q_v\)，按来源独立性分组后的可见支持数记为 \(m_v(C_t)\)。版本化 corroboration profile 为每个节点指定阈值 \(k_v\)，并定义：

\[
v\in V_t^{\mathrm{cov}}
\Longleftrightarrow
m_v(C_t)\geq k_v .
\]

其中，\(k_v=1\) 对应乐观 OR 策略，\(k_v=|Q_v|\) 对应保守 AND 策略，中间取值表示部分 corroboration。该推广不改写已有结果：C11 冻结 AND 主结果与 OR 敏感性仍作为两个端点；新增扫描均通过独立 profile 记录。粒度梯度 G0--G3 仍作为实验性结构代理，而不是 actor accuracy。默认数值阈值保留为冻结参照，同时在节点覆盖集合 \(\{0.60,0.65,0.70,0.75,0.80\}\)、边覆盖集合 \(\{0.50,0.55,0.60,0.65\}\)、阶段数集合 \(\{2,3\}\) 组成的 40 个组合上报告稳健性。最终输出由案例级 \(\texttt{support\_ceiling}\) 截断。


\subsection{信息边界}

动作的公开目标必须在执行前由动作模板、通道能力或调查请求生成并版本化，不能在读取 \(R(a)\) 后反向构造。运行时采用递归 allowlist：配置、状态和动作仅暴露批准字段，并禁止访问 \(recoverable\_claim\_ids\)、\(required\_claim\_ids\)、mask、seed、运行 ID、实际通道状态、Oracle 路径和结果标签。AFA、Depth-2、XGBoost 和 Logistic 的正式入口均执行公开视图约束及隐藏结果不变性测试。规划器返回 \(\texttt{action\_id}\) 后，执行器才读取完整动作并实现恢复结果。该隔离属于软件接口约束，而非密码学访问控制。

若：

\[
I(a)=\operatorname{Cover}(R(a))
\]

且 \(I(a)\) 由 \(R(a)\) 确定，则 \(I(a)\) 已泄漏该动作的节点级恢复范围。编译审计因此将该等式作为高风险信号；但集合偶然相等本身不构成充分泄漏证据，关键在于生成过程是否读取隐藏结果。

图~\ref{fig:control-loop} 展示调查控制闭环与信息边界。规划器仅访问公开动作目标、成本、当前缺口、预算及历史反馈；执行器和 Oracle 隐藏实际恢复集合与实现通道状态。
\begin{figure*}[t]
  \centering
  \begin{tikzpicture}[
    font=\small,
    node distance=9mm and 11mm,
    box/.style={draw, rounded corners=2pt, minimum height=10mm, minimum width=30mm, align=center, fill=blue!4},
    hidden/.style={box, fill=red!5, draw=red!65!black},
    flow/.style={-{Latex[length=2.2mm]}, thick},
    feedback/.style={-{Latex[length=2.2mm]}, thick, dashed, teal!70!black}
  ]
    \node[box] (state) {公开调查状态\\证据缺口、预算、历史};
    \node[box, right=of state] (planner) {非 Oracle 规划器\\选择 \texttt{action\_id} 或 STOP};
    \node[hidden, right=of planner] (executor) {执行器\\读取完整动作并采集};
    \node[box, right=of executor] (evidence) {反馈与状态更新\\新增 claims 或零收益};
    \node[hidden, below=of executor] (oracle) {隐藏域\\$R(a)$、真实通道状态、结果标签};

    \draw[flow] (state) -- node[above]{公开视图} (planner);
    \draw[flow] (planner) -- node[above]{\texttt{action\_id}} (executor);
    \draw[flow] (executor) -- node[above]{采集结果} (evidence);
    \draw[flow, red!65!black] (oracle) -- (executor);
    \draw[feedback] (evidence.south) to[out=-110,in=-70] node[below]{可审计反馈} (state.south);

    \node[draw=blue!65!black, dashed, rounded corners, fit=(state)(planner), inner sep=5pt,
          label=above:{\bfseries 规划器可见域}] {};
    \node[draw=red!65!black, dashed, rounded corners, fit=(executor)(oracle), inner sep=5pt,
          label=above:{\bfseries 执行/Oracle 隐藏域}] {};
  \end{tikzpicture}
  \caption{调查控制闭环与信息边界。规划器只读取执行前公开且版本化的状态；隐藏恢复集合和真实通道状态仅在动作选定后由执行器访问。}
  \label{fig:control-loop}
\end{figure*}


\subsection{决策、反馈和停止}

系统在预算约束下选择动作，执行后接收新增 claims 或零收益反馈。成功表示内部目标粒度在预算内达到，而不是 actor/campaign 正确率。设 STOP 的终端效用为 \(U(s)\)，则在有限时域内，STOP 最优当且仅当：

\[
U(s)\geq
\max_{a\neq \mathrm{STOP}}
\{-c_p(a)+\mathbb{E}[V^\star(s')|s,a]\}.
\]

单步增益小于成本并不足以推出停止，因为低即时收益动作可能解锁后续关键动作。若目标结构不可达、预算不足或 continuation value 不高于 STOP，则系统返回当前支持粒度、未覆盖关键节点以及停止原因。


\section{参数治理与实现}


\subsection{成本 profile}

运行器支持四种成本口径：

\begin{table*}[t]
  \centering
  \caption{成本 profile 的语义、状态与解释边界}
  \label{tab:cost-profiles}
  \resizebox{\textwidth}{!}{%
\begin{tabular}{llll}
\toprule
Profile & 含义 & 当前状态 & 可作何种解释 \\
\midrule
legacy & 原案例中的历史相对成本 & 正式运行完成；与旧输出逐字节兼容 & 冻结历史基线，不代表真实运营成本 \\
uniform & 所有非 STOP 动作为 1 & 正式运行完成 & 动作步数对照，去除人工成本排序 \\
rubric & E 努力、V 易失性、D 数据量、A 访问合规、R 侵入风险五分量 & 360 个独立人工评分分量待完成 & 冻结前不得进入正式实验 \\
measured & 动作级运营测量归一化后的连续成本 & 72/72 动作待测；基础设施已完成 & 当前不得报告数值或替代真实成本 \\
\bottomrule
\end{tabular}
  }
\end{table*}

每个正式 profile 均包含 ID、语义版本、状态和 SHA-256。rubric/measured 草案、动作覆盖不完整的 profile，以及非空旧结果目录均被拒绝。当前仅有 \(41/72\) 个动作具有案例级扫描总量，\(40/72\) 个动作具有产出量或事件跨度；动作级扫描字节、保留期和主机数均为：

\[
0/72 .
\]

案例总量不能冒充动作扫描量，事件跨度不能冒充保留期，旧成本也不能用于回填评分缺失。


\subsection{收益权重与先验 profile}

透明 M2 被解释为：

\[
VoI(a,s)-\alpha c_p(a)
\]

的领域线性化，而非具有唯一理论系数的真值模型。冻结实现保留原权重，新增实验仅通过外部 profile 扫描：

\[
\alpha\in
\{0,0.25,0.50,0.75,1.00,1.50,2.00\}.
\]

expected effects 和通道可靠性同时区分 measured、standard 与 expert-prior 来源；先验缩放仅改变规划器信念，而不改变执行器实际可靠性，从而避免将环境退化误写为 prior sensitivity。
\subsection{规划器}
比较对象包括透明 M2、M3a gap compatibility、XGBoost、Logistic/M3b、AFA-VOI Myopic、AFA-VOI Rollout-H3、Depth-2 和仅供上界解释的 Oracle。XGBoost 与 Logistic 只使用公开状态—动作特征；AFA 和 Depth-2 使用相同的公开通道先验；Oracle 可以读取隐藏结果但不属于部署候选。所有正式输出记录 profile 哈希、运行契约和输入哈希。

\section{实验设计}
\subsection{案例}
C01–C06 仅用于开发、训练或规则校验，不进入本文的独立评估案例数。C07–C10 是 DARPA TC E5 与 OpTC 的四个参数锁定 G3 案例；C11 是一个 OTRF APT29 仿真链，因冻结锚点自然缺失而截断为 G2；C12 是从生产 SOC 衍生数据经元数据和事件来源两级 Gate 筛出的一个双流 incident，支持上限为 G1。C11 的多 provider claims 来自同一主机归档，不等于独立网络传感器 corroboration；C12 的厂商 GraphML、自动 ATT\&CK 标签和 Disrupted 状态只作上下文，不作 actor truth。

\subsection{重复测量}
每个案例运行 45 个 mask/intensity/seed 条件，共 270 条重复条件。所有条件在 profile 对照中逐行配对。描述性汇总可以按 270 条条件计算成功次数和平均成本，但科学外推的独立 n 始终为 6。由于案例数很小且来源异质，本文不报告把重复条件当独立样本得到的置信区间或 p 值。

\subsection{终点}
主要内部终点为是否达到案例目标、成功条件下累计成本、相对 Oracle regret、零收益动作、过早停止和粒度越界。附加证据受限终点包括 over-attribution、abstention/降级和来源覆盖。actor/campaign accuracy 与分析师效用只有在外部 ground-truth bundle 完整时才计算；当前二者均不可识别。

\subsection{治理扫描}
除 legacy/uniform 对照外，本文报告 40 个阈值组合、6 个 claim/source-group k-of-n 配置、7 个 M2 $\alpha$ 值、expected-effect 专家先验缩放和 AFA/Depth-2 通道先验缩放。所有变体在查看结果前由 profile 固定，并写入新结果目录；不改写原始 \texttt{acquisition\_actions.json} 或旧结果。

\subsection{复现与完整性}
15 个正式输出由统一审计重新验证。测试状态为 411 passed、6 skipped、0 failed。该数字是软件验证记录，不是科学样本量。审计清单明确 \texttt{all\_experiments\_complete=false}：自动化实现已完成，但人工、测量和外部真值 Gate 仍开放。

\section{结果}
\subsection{冻结 legacy 结果}
在 C07–C10 的 180 个案例内重复条件中，M2 全部达到内部 G3 目标，平均成功成本为 4.5333；XGBoost、M3a、AFA-VOI 和公开 Depth-2 均未降低该聚合成本。C11 中，M2 为 45/45、平均成本 3.6667；XGBoost 和 Logistic 均为 3.0667，Depth-2 则在 1 个条件出现成功退化。C12 中，所有实现均保持零粒度越界，但策略排序再次变化，Depth-2 在该单一 G1 案例中与 Oracle 同为 0.8889。三层结果表明接口可以在不同上限下运行，却不支持固定策略排序跨来源泛化。

在统一的 C07–C12 ML 正式入口中：

\begin{table*}[t]
  \centering
  \caption{冻结 legacy 成本下的 ML 正式入口结果}
  \label{tab:legacy-results}
  \resizebox{\textwidth}{!}{%
\begin{tabular}{lllll}
\toprule
Planner & 达标条件 & Success & Mean successful cost & 解释 \\
\midrule
XGBoost & 270/270 & 1.0000 & 3.8926 & 冻结训练、公开特征 \\
M2 & 270/270 & 1.0000 & 3.8704 & 透明部署锚点 \\
Logistic/M3b & 267/270 & 0.9889 & 4.0000 & 3 个失败均保留 \\
\bottomrule
\end{tabular}
  }
\end{table*}

这些数字使用历史相对成本，只能回答冻结实现下发生了什么，不能回答真实运营中哪个规划器更便宜。

\subsection{Uniform 成本对照}
Uniform profile 把每个非 STOP 动作成本固定为 1，因此平均成本等于平均动作数。三个正式入口得到：

\begin{table*}[t]
  \centering
  \caption{单位成本口径下各正式入口的平均动作数}
  \label{tab:uniform-results}
  \resizebox{\textwidth}{!}{%
\begin{tabular}{llll}
\toprule
正式入口 & Planner & 达标条件 & Mean uniform cost \\
\midrule
AFA & M2 & 270/270 & 1.7296 \\
AFA & AFA-VOI Myopic & 270/270 & 1.7333 \\
AFA & AFA-VOI Rollout-H3 & 270/270 & 1.7778 \\
AFA & Oracle & 270/270 & 1.4852 \\
Depth-2 & M2 & 270/270 & 1.7296 \\
Depth-2 & Depth-2 & 270/270 & 1.8963 \\
Depth-2 & Oracle & 270/270 & 1.4852 \\
ML & XGBoost & 270/270 & 1.6852 \\
ML & Logistic/M3b & 270/270 & 1.6630 \\
ML & M2 & 270/270 & 1.7296 \\
ML & M3a & 270/270 & 1.7333 \\
ML & Oracle & 270/270 & 1.4852 \\
\bottomrule
\end{tabular}
  }
\end{table*}

成功闭环在 uniform 下没有退化，但 legacy 中 M2 略低于 XGBoost，uniform 中 XGBoost 与 Logistic/M3b 都低于 M2。该排序变化是成本口径敏感性的直接证据。Uniform 只测动作数；它既不证明学习器真实更省资源，也不能替代 rubric 或 measured cost。

\subsection{阈值、corroboration 与 \texorpdfstring{$\alpha$}{alpha} 稳健性}

\begin{table*}[t]
  \centering
  \caption{参数治理扫描与动作序列稳定性}
  \label{tab:governance-scan}
  \resizebox{\textwidth}{!}{%
\begin{tabular}{lllll}
\toprule
治理项 & 变体数 & M2 success 范围 & 最低动作序列一致率 & 最大粒度越界率 \\
\midrule
Cost（legacy/uniform） & 2 & 1.0000–1.0000 & 0.9111 & 0 \\
G3 阈值网格 & 40 & 1.0000–1.0000 & 1.0000 & 0 \\
k-of-n corroboration & 6 & 1.0000–1.0000 & 0.8148 & 0 \\
M2 $\alpha$ & 7 & 1.0000–1.0000 & 0.8407 & 0 \\
\bottomrule
\end{tabular}
  }
\end{table*}
因此可以说 M2 的内部达标与不越界在这些扫描中稳定；不能说动作选择或成本排序不受参数影响。阈值网格在当前六例上没有触发序列变化，主要因为冻结关键节点和支持上限主导达标；这不是阈值已获外部校准的证据。

\subsection{expected effects 与 planner prior}
修正后的 W6 实验把 planner belief 与执行可靠性分离。相对 legacy，dev-measured base 在 C11 的 3 个重复条件使 M2 过早停止；将专家先验提高 25\% 修复这 3 个条件。仅改变 M2 不消费的通道信念字段不会改变结果，这一零效应是契约行为而不是先验稳健性证明。

对实际消费通道 prior 的规划器，0.75 倍缩放使 AFA-VOI Myopic 在 C09 出现 2 个成功损失；Rollout-H3 的成功数不变，但动作序列在部分条件发生变化。Depth-2 出现 1 个成功损失和 1 个成功修复，聚合 success 相互抵消。1.25 倍缩放未改变 AFA 或 Depth-2 的动作和结果。M2 与 Oracle 不读取该 planner prior，因而保持不变。结论是低估通道可靠性会实质改变部分策略，而不是某个固定 prior 已被证实正确。

\subsection{外部终点和人工构念 Gate}
外部 ground-truth 接口已实现，但 12 个案例中当前没有可接受的 actor/campaign 真值或分析师效用记录；\texttt{external\_actor\_accuracy} 因而为 \texttt{not\_identifiable}。内部 success、G1–G3 或厂商标签均不能填补该空缺。

Round 1 中，27 个 Claim 项和 27 个 Intent 项具备来源独立的双人可比记录。Claim 支持度加权 $\kappa=-0.1455$；Intent mean Jaccard=0.3673、micro-F1=0.4878，均未过预注册门槛。另 60 个粒度项目的 A/B 源文件 SHA-256 完全相同，故全部作废。第三人裁决不能把首轮一致性改写为通过；Round 2 尚未启动。Claim、Intent 和粒度因此仍是待验证工程构念。

\section{讨论}
\subsection{可以成立的主张}
第一，Project05 已实现一个由来源声明、公开动作视图、通道执行、反馈、STOP 和支持上限组成的调查控制闭环。第二，运行时 allowlist 和结果不变性测试使非 Oracle 规划器不能通过主入口读取恢复集合、mask、实际通道状态或标签。第三，六个来源异质的案例表明闭环可在 G3、G2 和 G1 上限下运行，并能保留自然降级。第四，参数治理将成本、阈值、corroboration、收益权重和先验从隐含常量转化为版本化、可审计和可配对比较的实验对象。

\subsection{不能成立的主张}
本文没有 actor/campaign 正确性或分析师效用真值，故不能声称提高归因准确率。独立案例只有 6 个，不能声称跨组织泛化。AFA 与非短视方法是领域适配，不能称为 NOCTA、AFABench 或 WinRegRL 的官方同任务复现。M2 在多组扫描中保持成功不等于全局最优；uniform 下的排序变化和 prior 敏感性明确否定了这种解释。Rubric 和 measured profile 尚未冻结，不能把 legacy 或 uniform 称为真实运营成本。

\subsection{参数治理改变了什么}
参数治理没有“找到更漂亮的一组数”，也没有回写旧结果。它改变的是证据标准：每个数必须有来源类型、版本和哈希；影响 planner-visible 字段的 profile 必须先冻结再运行；比较必须逐条件配对；结果方向若只在单一口径成立，则只能报告为条件性结果。由此，论文的主要经验结论从“M2 成本最低”收缩为“闭环成功较稳定，但动作路径和成本排序依赖尚未完全外部校准的参数”。这一较弱结论更符合现有证据。

\subsection{LLM 和学习器的位置}
LLM 可以用于 schema 约束的离线 claim 编译、来源解释和异常提示，学习器可以在公开特征上估计动作价值。但在来源真实性、构念复现和外部终点未关闭前，它们都不应获得不可审计的在线权限。当前 XGBoost 和 Logistic 结果说明学习器可以接入统一契约；uniform 排序又说明其表现必须与成本口径共同解释，而不能凭模型类型获得普遍优势。

\section{局限性与开放 Gate}

\begin{enumerate}[leftmargin=*,label=(\arabic*)]
  \item 独立案例数为 6；270 个重复条件不能增加外部效度。
  \item 外部 actor/campaign GT 与分析师效用记录为 0，真实归因能力不可识别。
  \item Rubric 尚有 360 个独立评分分量待两名真实评分者完成；Round 2 尚未启动。
  \item Measured profile 的 72 个动作均缺正式运营测量，尤其动作级字节、保留期和主机数为 0/72。
  \item G0–G3 是结构代理；40 个阈值组合的稳定性不等于与人类结论一致。
  \item C11 只是一个仿真链，C12 只是一个 G1 incident；各自的策略排序不能单独外推。
  \item k-of-n profile 统一了 OR/AND 语义，但来源独立性的判定本身仍需领域和人工复核。
  \item DP 只作完整信息上界；合成环境中的非短视必要性不估计真实调查中的发生率。
\end{enumerate}

\section{结论}
本文把不完整证据下的 APT 前置调查定义为信息边界约束的序贯控制问题：规划器依据公开且版本化的证据缺口状态、动作目标、成本和先验选择采集或 STOP，执行器随后返回新增证据或零收益反馈，最终输出由来源和支持上限约束。六个独立案例上的 270 个案例内重复条件表明，该闭环在 legacy 与 uniform 口径下都能维持较高内部成功；同时，成本排序、动作路径和部分规划器的成功会随成本与 prior 变化。

\section*{数据与代码可用性}
原始 DARPA、OTRF 和生产 SOC 衍生数据受各自发布与访问条件约束，仓库不重新分发受限大体量原始数据。案例配置、来源指针、profile schema、profile ID/版本/SHA-256、运行契约、正式结果、配对审计和回归测试保存在 GitHub 仓库。统一审计清单位于 09-experiments/results/nonhuman\_completion\_audit\_v0.1/。对任何结果复现，必须同时报告独立案例数、重复条件定义、成本 profile 和运行契约；不得只复制汇总 CSV 而省略 profile 哈希。

\balance
\begin{thebibliography}{99}
\bibitem{ref1} Zhiqiang Xu; Pengcheng Fang; Changlin Liu; Xiangyu Xiao; Yinqian Wen; Dinghao Meng. (2022). DEPCOMM: graph summarization on system audit logs for attack investigation. doi:\url{https://doi.org/10.1109/SP46214.2022.9833632}

\bibitem{ref2} Fan Yang; Jiacen Xu; Chunlin Xiong; Zhou Li; Kehuan Zhang. (2023). PROGRAPHER: an anomaly detection system based on provenance graph embedding. USENIX Security 2023. \url{https://www.usenix.org/conference/usenixsecurity23/presentation/yang-fan}

\bibitem{ref3} Zijun Cheng; Qiujian Lv; Jinyuan Liang; Yan Wang; Degang Sun; Thomas Pasquier; Xueyuan Han. (2024). Kairos: practical intrusion detection and investigation using whole-system provenance. doi:\url{https://doi.org/10.1109/SP54263.2024.00005}

\bibitem{ref4} Jia, Zian; Xiong, Yun; Nan, Yuhong; Zhang, Yao; Zhao, Jinjing; Wen, Mi. (2024). MAGIC: Detecting Advanced Persistent Threats via Masked Graph Representation Learning. 33rd USENIX Security Symposium (USENIX Security 24), pp. 5197–5214. \url{https://www.usenix.org/conference/usenixsecurity24/presentation/jia-zian}

\bibitem{ref5} Milajerdi, Sadegh M.; Eshete, Birhanu; Gjomemo, Rigel; Venkatakrishnan, V. N.. (2019). POIROT: Aligning Attack Behavior with Kernel Audit Records for Cyber Threat Hunting. doi:\url{https://doi.org/10.1145/3319535.3363217}

\bibitem{ref6} Kiavash Satvat; Rigel Gjomemo; V. N. Venkatakrishnan. (2021). EXTRACTOR: extracting attack behavior from threat reports. doi:\url{https://doi.org/10.1109/EuroSP51992.2021.00046}

\bibitem{ref7} Li, Jingwen; Zhang, Ru; Liu, Jianyi; Zhao, Wanguo. (2025). CLIProv: A Contrastive Log-to-Intelligence Multimodal Approach for Threat Detection and Provenance Analysis. \url{https://arxiv.org/abs/2507.09133}

\bibitem{ref8} Qiu, Xuebo; Lv, Mingqi; Zhang, Yimei; Chen, Tieming; Zhu, Tiantian; Song, Qijie; Ji, Shouling. (2025). APT-CGLP: Advanced Persistent Threat Hunting via Contrastive Graph-Language Pre-Training. \url{https://arxiv.org/abs/2511.20290}

\bibitem{ref9} Emanuele Mezzi; Fabio Massacci; Katja Tuma. (2025). Large language models are unreliable for cyber threat intelligence. arXiv. \url{https://arxiv.org/abs/2503.23175}

\bibitem{ref10} Max Van Der Horst; Ricky Kho; Olga Gadyatskaya; Michel Mollema; Michel Van Eeten; Yury Zhauniarovich. (2025). High stakes, low certainty: evaluating the efficacy of high-level indicators of compromise in ransomware attribution. USENIX Security 2025. \url{https://www.usenix.org/conference/usenixsecurity25/presentation/van-der-horst}

\bibitem{ref11} Aronsson, Linus; Rahbar, Arman; Chehreghani, Morteza Haghir. (2025). A Survey on Active Feature Acquisition Strategies. \url{https://arxiv.org/abs/2502.11067}

\bibitem{ref12} Valancius, Michael; Lennon, Maxwell; Oliva, Junier. (2024). Acquisition Conditioned Oracle for Nongreedy Active Feature Acquisition. Proceedings of the 41st International Conference on Machine Learning, 235, pp. 48957–48975. \url{https://proceedings.mlr.press/v235/valancius24a.html}

\bibitem{ref13} Dinh, Dzung; Chen, Boqi; Qu, Yunni; Niethammer, Marc; Oliva, Junier. (2025). NOCTA: Non-Greedy Objective Cost-Tradeoff Acquisition for Longitudinal Data. \url{https://arxiv.org/abs/2507.12412}

\bibitem{ref14} Sch\"{u}tz, Valter; Wu, Han; Rezvan, Reza; Aronsson, Linus; Haghir Chehreghani, Morteza. (2026). AFABench: A Generic Framework for Benchmarking Active Feature Acquisition. Proceedings of the 32nd ACM SIGKDD Conference on Knowledge Discovery and Data Mining. doi:\url{https://doi.org/10.1145/3770855.3817493}

\bibitem{ref15} Ghanem, Mohamed Chahine; Wojtczak, Dominik; Benkhelifa, Elhadj; Kheddar, Hamza; Nepomuceno, Erivelton G.; Li, Wanpeng. (2026). Leveraging Reinforcement Learning for an Efficient Windows Registry Analysis during Cyber Incident Response. Scientific Reports. doi:\url{https://doi.org/10.1038/s41598-026-57787-6}

\bibitem{ref16} Turney, Peter D.. (2000). Types of Cost in Inductive Concept Learning. Proceedings of the Workshop on Cost-Sensitive Learning at ICML 2000, pp. 15–21.

\bibitem{ref17} Saar-Tsechansky, Maytal; Melville, Prem; Provost, Foster. (2009). Active Feature-Value Acquisition. Management Science, 55(4), pp. 664–684. doi:\url{https://doi.org/10.1287/mnsc.1080.0952}

\bibitem{ref18} Howard, Ronald A.. (1966). Information Value Theory. IEEE Transactions on Systems Science and Cybernetics, 2(1), pp. 22–26. doi:\url{https://doi.org/10.1109/TSSC.1966.300074}

\bibitem{ref19} Brezinski, Dominique; Killalea, Tom. (2002). Guidelines for Evidence Collection and Archiving. Internet Engineering Task Force. doi:\url{https://doi.org/10.17487/RFC3227}

\bibitem{ref20} Kent, Karen; Chevalier, Suzanne; Grance, Tim; Dang, Hung. (2006). Guide to Integrating Forensic Techniques into Incident Response. National Institute of Standards and Technology. doi:\url{https://doi.org/10.6028/NIST.SP.800-86}
\end{thebibliography}
\end{document}